\documentclass[../main.tex]{subfiles}
\begin{document}

This appendix provides additional detailed results supplementing Chapter~\ref{ch:results}.

\section{Population Structure from Core SNPs}

Table~\ref{tab:popstruct} shows the number of lineages and the size of the largest lineage at different SNP-distance thresholds, evidencing that the dataset is almost single-clonal.

\begin{table}[H]
\centering
\caption{Population structure (core SNP clustering).}
\label{tab:popstruct}
\begin{tabular}{ccc}
\hline
SNP-distance threshold & No. of lineages & Largest lineage (\%) \\ \hline
0.05 & 14 & 97.4 \\
0.10 & 13 & 97.5 \\
0.20 & 13 & 97.5 \\
0.30 & 8  & 99.1 \\ \hline
\end{tabular}
\end{table}

\section{QRDR Mutation Association with Resistance}

Table~\ref{tab:qrdr} lists the resistance rate by mutation status at the QRDR positions for CIP and NAL.

\begin{table}[H]
\centering
\caption{QRDR mutations and resistance rate (R-rate).}
\label{tab:qrdr}
\begin{tabular}{llcc}
\hline
Drug & Position & R-rate if mutated & R-rate if wild-type \\ \hline
CIP & gyrA83 & 0.76 & 0.04 \\
CIP & gyrA87 & 0.87 & 0.13 \\
CIP & parC80 & 0.88 & 0.15 \\
NAL & gyrA83 & 0.77 & 0.06 \\
NAL & gyrA87 & 0.84 & 0.04 \\
NAL & parC80 & 0.75 & 0.08 \\ \hline
\end{tabular}
\end{table}

\end{document}
